The proposed DSL is just a first step towards a more diverse,  
transparent and agentic collaboration in software development. 
Open challenges and future directions are discussed below.

\smallskip
\noindent\textbf{Usability improvement via new syntaxes.}
To facilitate the adoption of our DSL by less technical users, we envision offering alternative UIs, like chatbots, to enhance usability.
%  to define the policies, enhancing usability through intuitive UIs or chatbot interfaces. 
This will be especially useful when introducing our DSL to manage governance in other, non-software, domains.
% Furthermore, we could leverage some technologies, such as BESSER Agentic Framework\footnote{\url{https://github.com/BESSER-PEARL/BESSER-Agentic-Framework}}, to create a chatbot interface that guides users through the policy definition process.
%This would broaden accessibility beyond technical users, fostering more inclusive governance design.

\smallskip
\noindent\textbf{Longitudinal study.}
To better understand if explicit governance rules help to increase 
%the number of contributions in an OSS project,
OSS contributions,
% (and, if so, what types of policies are more helpful to that), 
we plan a longitudinal study where we could compare the evolution of contributions in projects before and after they committed to an explicit governance model. 
This same study could be useful to see whether projects are sticking to such governance model or they just ignore it even after committing to it. 

%ly, we intend to investigate usability aspects and gather feedback from practitioners to refine the language and its supporting tools.

\smallskip
\noindent\textbf{Improved impact of Multi-agent systems in software development.}
We plan to test whether our governance layer enables more complex multi-agent systems to collaborate, in a controlled way, to solve more ambitious software engineering tasks. 
By leveraging SWE-bench~\cite{DBLP:conf/iclr/JimenezYWYPPN24}, we can systematically evaluate how governance policies affect agent collaboration, task completion rates, and adherence to defined constraints.
% This evaluation will provide insights into how formal governance structures impact autonomous agent behavior and overall system performance in software engineering contexts.

\smallskip
\noindent\textbf{Participant identification and bias mitigation.}
A problem some projects face is identifying the right people to play certain roles. 
Given a governance model, it would be possible to analyze past project data to propose candidates to fill the roles required in the governance model. 
% These candidates could then be validated by the community. 
A real-time monitoring component should also be helpful to proactively change the candidates to continuously improve the diversity of the collaborations. Indeed, a key element of this (semi)automatic participant analysis could be devoted to make sure the participants list is not biased against certain communities and that it respects the profile requirements. 
%This would enable automated role adjustments, proactive governance violation detection, and adaptive policy enforcement incorporating positive and negative discrimination strategies to maintain fair representation.

%\smallskip
%\noindent\textbf{Bias mitigation and diversity enhancement.}
%Finally, we plan to introduce mechanisms that actively promote diversity in governance decisions and team composition.
%This includes automatic role identification in governance processes and implementing diversity metrics that can be continuously monitored and enforced, thus maintaining inclusive participation across different profile dimensions.

\smallskip
\noindent\textbf{Governance beyond software.}
We believe other domains could also benefit from an explicit and precise definition of their governance policies based on our DSL. 
% This broader applicability underscores the potential impact of a generalized governance framework for mixed-initiative systems and is a key direction for future exploration.
For instance, making explicit the governance systems of organizations (such as NGOs) would enable systematic comparison of their decision-making structures and their assessment of inclusivity. % across different organizational contexts.